\bookchapter{Capítulo 22: Como a IA funciona no celular}

\chapterintro{Bem-vindo ao “show de mágica” mais fascinante do seu celular, com a estrela principal: a Inteligência Artificial. As fotos incríveis que você tira “com um clique”, o assistente de voz que responde prontamente, o teclado que parece te conhecer cada vez melhor — por trás de tudo isso está esse cérebro invisível.}

\booksection{Câmera com IA: seu “laboratório de revelação digital” no celular}

A fotografia computacional com IA de hoje é como ter uma equipe de editores profissionais + fotógrafos embutida diretamente no botão do obturador. Antes mesmo de você pressionar o botão, a IA já está trabalhando: escaneando a cena, identificando “isto é céu azul”, “isto é um rosto”, “isto é cena noturna”.

\booksection{Assistentes de voz: como seu “secretário de voz” te entende?}

\begin{enumerate}
  \item \textbf{Ouvir (reconhecimento de fala)}: converte ondas sonoras em texto.
  \item \textbf{Entender (compreensão de linguagem natural)}: compreende a intenção, extrai entidades-chave.
  \item \textbf{Executar e responder}: aciona o app ou serviço correspondente, depois sintetiza a voz para "falar" com você.
\end{enumerate}

\booksection{IA no dispositivo (on-device) vs. IA na nuvem}

A IA no dispositivo mantém o processamento dentro do celular: rápida, mais privada, funciona offline. A IA na nuvem precisa de internet, mas pode lidar com tarefas mais complexas.

\booksection{Grandes modelos de IA e o celular: seu assistente inteligente de bolso}

Em 2025-2026, a explosão dos grandes modelos de IA (como DeepSeek R1, ChatGPT-5, Wenxin Yiyan 5.0, Tongyi Qianwen 3, entre outros) transformou profundamente a forma como usamos o celular. O telefone deixou de ser apenas uma ferramenta que executa comandos e se tornou um parceiro inteligente capaz de compreender, raciocinar e criar. Os chips de IA no dispositivo (NPU com poder de computação superior a 80 TOPS) permitiram, pela primeira vez, a inferência de grandes modelos diretamente no celular, sem precisar de internet.

\booksubsection{O que é um grande modelo de IA?}

Grandes modelos de IA (Large Language Model, LLM) são sistemas de IA treinados com enormes quantidades de dados textuais, capazes de compreender e gerar linguagem natural. Imagine-o como um “supergênio” que leu praticamente todo o texto da internet e consegue responder perguntas, escrever artigos, traduzir, programar e analisar dados.

\booksubsection{Aplicações de grandes modelos de IA no celular:}

\begin{itemize}
  \item \textbf{Assistente de conversa}: apps como DeepSeek, ChatGPT, Doubao, Kimi, que permitem perguntar, traduzir, redigir textos, resumir artigos a qualquer momento. Mais inteligentes que os mecanismos de busca tradicionais — não te dão apenas links, mas respostas diretas.
  \item \textbf{Escrita e criação}: deixe a IA escrever e-mails, resumir atas de reunião, polir textos, gerar textos criativos.
  \item \textbf{Geração e edição de imagens}: IA remove pessoas indesejadas de fotos com um clique, expande bordas de imagens, converte estilos e até gera imagens a partir de descrições textuais.
  \item \textbf{Tradução em tempo real}: tradução em tempo real de chamadas, tradução de conversas presenciais, tradução por câmera — viajar para o exterior sem barreiras linguísticas.
  \item \textbf{Apoio ao aprendizado}: a IA pode explicar conceitos, criar exercícios, corrigir tarefas — um tutor particular de bolso.
\end{itemize}

\booksubsection{Grandes modelos no dispositivo: o celular rodando IA localmente}

Os smartphones topo de linha mais recentes já conseguem executar pequenos grandes modelos de IA localmente (geralmente com cerca de 7 bilhões de parâmetros), trazendo as seguintes mudanças:

\begin{itemize}
  \item \textbf{Sem necessidade de internet}: funções de IA funcionam offline, no avião, no metrô.
  \item \textbf{Privacidade mais segura}: suas conversas e dados não são enviados para a nuvem, tudo é processado localmente no celular.
  \item \textbf{Resposta mais rápida}: sem latência de rede, a resposta da IA é quase instantânea.
  \item \textbf{Limitações atuais}: os modelos no dispositivo ainda não alcançam a capacidade dos grandes modelos na nuvem. Tarefas complexas ainda exigem internet. Mas com o aumento do poder de computação dos chips de IA, os modelos locais ficarão cada vez mais fortes.
\end{itemize}

\booksubsection{Integração de IA a nível de sistema:}

\begin{itemize}
  \item \textbf{Apple Intelligence}: sistema de IA introduzido no iOS 18, integrando ferramentas de escrita, varinha mágica de imagem, Siri aprimorada, resumos de notificações e mais. Prioriza processamento no dispositivo, recorrendo à nuvem apenas para tarefas complexas.
  \item \textbf{Modelo Pangu da Huawei}: integrado ao sistema HarmonyOS, a assistente Xiaoyi ganha capacidades de grande modelo, suportando resumo de documentos, reconhecimento inteligente de tela, edição de fotos com IA, etc.
  \item \textbf{Xiaomi Pengpai AI}: capacidades de IA a nível de sistema, suportando busca de fotos com IA, resumos com IA, legendas em tempo real com IA, conversa com grandes modelos no dispositivo, etc.
  \item \textbf{Samsung Galaxy AI}: funções como tradução em tempo real de chamadas, resumos de notas com IA, geração e edição de imagens, busca com seleção circular, etc.
\end{itemize}

\booksection{Cenários típicos de aplicação e prática}

Para você sentir de forma mais intuitiva os “superpoderes” da IA no celular, vamos usar cinco cenários específicos para mostrar como a IA “salva o dia” no cotidiano:

\booksubsection{Cenário 1: Tradução em tempo real — o intérprete de viagem internacional}

\textbf{Desempenho real}: pedindo comida em Paris sem falar francês? O app “Traduzir” do iPhone suporta modo de conversa presencial para duas pessoas, cada um vendo legendas no seu próprio idioma, podendo até fazer tradução simultânea. No lado Android, a tradução em tempo real de chamadas do Samsung Galaxy AI permite que você fale em chinês ao telefone enquanto a outra pessoa ouve em inglês, e vice-versa.

\textbf{Princípio por trás}: pipeline de três etapas — reconhecimento de fala (ASR) tradução automática (NMT) síntese de fala (TTS). A latência de ponta a ponta já pode ser controlada em menos de 1 segundo.

\booksubsection{Cenário 2: Remoção com IA — salvando fotos de turista “arruinadas” por pedestres}

\textbf{Desempenho real}: tirou uma selfie incrível num ponto turístico, mas o fundo está cheio de outros turistas? Use a “remoção mágica” para circular os pedestres — o celular os identifica automaticamente, os “apaga” e preenche o fundo de forma inteligente. O resultado muitas vezes é impressionante.

\textbf{Princípio por trás}: o modelo de segmentação de imagem (Inpainting) primeiro identifica e remove a área selecionada, depois “imagina” um preenchimento razoável com base nos pixels ao redor e na compreensão semântica. O modelo de IA por trás disso foi treinado com milhões de imagens.

\booksubsection{Cenário 3: Edição de fotos com IA — torne-se um “mestre da fotografia” com um clique}

\textbf{Desempenho real}: tirou uma foto de paisagem cinzenta em dia nublado? O “céu mágico” da IA pode trocar o céu para céu azul, pôr do sol ou até céu estrelado com um clique. A expansão com IA pode transformar uma composição vertical em horizontal, com o celular “imaginando” o cenário fora do quadro.

\textbf{Princípio por trás}: baseado em modelos de difusão (Diffusion Model) da IA generativa (Generative AI), que compreendem o conteúdo semântico da imagem e geram novos pixels. Esta é a “fronteira mais avançada” do processamento de imagens com IA atual.

\booksubsection{Cenário 4: Processamento de documentos — o multiplicador de eficiência para profissionais}

\textbf{Desempenho real}: após uma reunião, entregue a gravação de áudio para a IA. Ela organiza automaticamente uma ata estruturada, extrai itens de ação e gera um resumo. Um relatório PDF de 10 páginas em inglês pode ser resumido pela IA em 10 segundos, extraindo as conclusões principais.

\textbf{Princípio por trás}: grandes modelos de contexto longo (Long Context LLM) podem “ler” documentos com dezenas de milhares de palavras, realizando múltiplas rodadas de raciocínio e extração de informações. Apple Intelligence, Huawei Xiaoyi, WPS AI, Tencent Yuanbao, entre outros, possuem essa capacidade.

\booksubsection{Cenário 5: Busca de imagens e reconhecimento de objetos com IA — seu “mecanismo de busca visual”}

\textbf{Desempenho real}: viu uma flor desconhecida na rua? Abra a câmera e a IA imediatamente te diz “é um hibisco”. Procurando no álbum aquela foto tirada na praia no ano passado? Pesquise “mar” e a IA automaticamente encontra todas as fotos com mar, sem precisar rolar manualmente.

\textbf{Princípio por trás}: modelos de visão computacional (CV) vetorizam (Embedding) o conteúdo das imagens e depois fazem correspondência de similaridade com consultas textuais. Essa “busca multimodal” é um grande avanço no gerenciamento de álbuns de fotos no celular.

\booksection{Otimização de prompts: a “fórmula de comunicação” para fazer a IA te obedecer}

A qualidade da resposta da IA depende muito de como você “pergunta”. A mesma pergunta, formulada de maneiras diferentes, pode gerar respostas completamente diferentes da IA. A seguir, apresentamos uma abordagem de comunicação\textbf{das mais práticas e fáceis de usar: “Papel-Contexto-Tarefa-Restrições”}, com muitos exemplos.

\begin{tcbinfo}{Por que preciso de uma estrutura?}
Muitas pessoas usam IA como usam o Baidu — digitam algumas palavras-chave. Mas a IA não é um mecanismo de busca; é um modelo de linguagem que “lê, compreende e depois responde”.\textbf{Dê a ela mais contexto, diga claramente o que você quer, limite-a para não se desviar}, e ela te dará resultados precisos. A estrutura é a tradução do que “uma boa resposta” significa na sua cabeça para instruções que a IA entende.

\end{tcbinfo}

\booksubsection{Papel (Role): diga à IA quem ela é}

Dê à IA uma “identidade” específica, e ela usará o conhecimento especializado e o estilo de expressão da área correspondente.\textbf{Quanto mais específico for o papel, mais profissional será o resultado}.

\booksubsection{Exemplos práticos:}

\begin{itemize}
  \item "Você é\textbf{um engenheiro de reparo de celulares Xiaomi com 10 anos de experiência}." (problema técnico)
  \item "Você é\textbf{um professor de física do ensino médio que é ótimo em explicar conceitos complexos de forma simples}." (pergunta de divulgação científica)
  \item "Você é\textbf{um editor de livros com bom gosto estético, que lê muito e tem estilo próprio}." (problema de escrita)
  \item "Você é\textbf{um blogueiro de culinária que prepara café da manhã para trabalhadores todos os dias}." (pergunta do dia a dia)
  \item "Você é\textbf{um médico de medicina tradicional chinesa paciente, gentil e que nunca zomba de iniciantes}." (pergunta de saúde)
\end{itemize}

\textbf{Por que funciona?}O papel define o\textbf{escopo de conhecimento}, o\textbf{estilo de expressão}e a\textbf{profundidade profissional}da IA. Sem um papel, a IA é um generalista que “sabe de tudo mas não é especialista em nada”. Com um papel, ela “entra” naquela identidade e pensa como aquele profissional.

\booksubsection{Contexto (Background): diga à IA sua situação}

Explique antecipadamente “quem você é, sua situação atual, suas restrições”, e a IA dará uma resposta\textbf{realmente adequada para você}, em vez de respostas genéricas.

\booksubsection{Exemplos práticos:}

\begin{itemize}
  \item "Meu celular é um\textbf{Xiaomi 14, comprei faz 6 meses, está esquentando muito durante jogos}." (condição do dispositivo)
  \item "Sou\textbf{uma senhora aposentada de 60 anos, com visão ruim, não sei mexer bem em smartphones}." (perfil do usuário)
  \item "Estou\textbf{me preparando para o exame CET-6 de inglês em junho de 2026, nível intermediário}." (contexto de estudo)
  \item "Abri\textbf{uma pequena cafeteria em uma cidade de terceiro nível, faturamento mensal de 30 mil, dois funcionários}." (contexto de negócio)
  \item "\textbf{Acabei de ter um bebê, estou de resguardo em casa, dormindo muito pouco}." (fase da vida)
\end{itemize}

\textbf{Por que funciona?}As informações de contexto permitem que a IA saiba\textbf{“para quem está servindo”}e\textbf{“em quais condições”}. A mesma pergunta, explicada para um estudante e para um pai, é diferente; o conselho para alguém da cidade e para alguém do interior também é diferente. Sem contexto, a IA só pode dar a resposta do “mínimo denominador comum”, geralmente sem impacto.

\booksubsection{Tarefa (Task): diga claramente à IA o que fazer}

Esta etapa deve começar com\textbf{verbos}(escreva, liste, compare, traduza, explique, resuma...), descrevendo a tarefa de forma\textbf{específica e executável}. Tarefas vagas só geram respostas vagas.

\booksubsection{Exemplos práticos:}

\begin{itemize}
  \item "Por favor,\textbf{liste 5 permissões que devo desativar imediatamente}e explique o motivo." (ação clara)
  \item "Por favor,\textbf{traduza este e-mail em inglês para o português}, mantendo o tom comercial." (tarefa específica)
  \item "Por favor,\textbf{compare estes quatro celulares na faixa de 3000 reais}e faça uma tabela de prós e contras." (formato claro)
  \item "Por favor,\textbf{explique 'o que é 5G' de um jeito que uma criança do primário entenda}." (estilo claro)
  \item "Por favor,\textbf{escreva três textos para vídeos curtos da minha lanchonete de café da manhã}, direcionados para trabalhadores, casais e pais com crianças." (segmentação clara)
\end{itemize}

\textbf{Por que funciona?}A parte da tarefa é o “centro de comando” da ação da IA.\textbf{Quanto mais claros os verbos e mais específicos os requisitos, mais a IA sabe “o que entregar”}. Palavras vagas como “dá um jeito”, “faz aí”, “dá uma olhada” são o maior inimigo da IA — ela vai tentar adivinhar o que você quer, e o resultado geralmente se desvia.

\booksubsection{Restrições (Constraints): diga à IA o que NÃO fazer}

“Restrições” soa como “limitação”, mas na verdade é uma ferramenta poderosa para\textbf{fazer a IA convergir para o resultado mais adequado para você}. Dimensões comuns incluem: número de palavras, formato, tom, itens a evitar, itens obrigatórios.

\booksubsection{Exemplos práticos:}

\begin{itemize}
  \item "\textbf{No máximo 200 palavras, em tópicos}." (tamanho + formato)
  \item "\textbf{Tom acolhedor e natural, como uma conversa entre amigos}, sem usar linguagem burocrática como 'primeiramente', 'em segundo lugar', 'ademais'." (tom)
  \item "\textbf{Não use termos técnicos}; se precisar usar um termo especializado, explique-o na sequência." (restrição de vocabulário)
  \item "\textbf{Cada parágrafo com no máximo 3 linhas}, para facilitar a leitura no celular." (formatação)
  \item "\textbf{Deve incluir: cenário de uso, prós e contras, compatibilidade com meu orçamento}." (itens obrigatórios)
  \item "\textbf{Não dê recomendações específicas de investimento}, apenas explique os princípios." (limites)
\end{itemize}

\textbf{Por que funciona?}Sem restrições, a IA tende a “extrapolar” — ou fica muito acadêmica, ou muito longa, ou foge do tema. Restrições claras = “critérios de entrega” claros. Alinhar expectativas de ambos os lados torna o resultado controlável.

\booksubsection{Avaliação da qualidade do prompt: 5 características de um bom prompt}

Escrever um bom prompt não é algo que se faz de uma vez.\textbf{Antes de enviar, verifique estes cinco itens}:

\begin{enumerate}
  \item \textbf{Específico, não abstrato}: não é "escreva um artigo", mas "escreva um guia de compra de celular de 800 palavras, com 3 seções, voltado para mães de primeira viagem".
  \item \textbf{Inclui contexto, não é isolado}: diga à IA "quem sou eu, por que estou perguntando isso", para que ela dê uma resposta "personalizada para você".
  \item \textbf{Tem formato de saída claro}: é tabela, lista, parágrafo ou JSON? Especifique antecipadamente.
  \item \textbf{Restrições adequadas}: número de palavras, estilo, itens a evitar — escolher um desses já traz uma melhoria significativa.
  \item \textbf{Pode ser verificado pela IA}: após perguntar, peça para a IA repetir seus requisitos; se ela não consegue descrevê-los, é porque você não os descreveu bem.
\end{enumerate}

\booksubsection{Problemas comuns e soluções}

\begin{center}
\begin{tabularx}{\linewidth}{|X|X|}
\hline
\textbf{Sintoma} & \textbf{Solução} \tabularnewline
\hline
Resposta muito genérica, sem utilidade prática & Adicione contexto (quem você é / qual cenário) \tabularnewline
Resposta muito acadêmica / formal demais & Adicione restrição de tom + lista de linguagem "burocrática" a evitar \tabularnewline
Resposta fora do tema / não responde a pergunta & Reformule a tarefa com verbos e exija "responda apenas isto" \tabularnewline
Resposta muito longa, impossível de ler & Limite forçado de palavras (máx. 200 palavras) \tabularnewline
IA inventa coisas com convicção & Adicione "se não tiver certeza, marque como 'não tenho certeza'" + verificação independente dos fatos importantes \tabularnewline
Primeira versão insatisfatória & Não reescreva; acrescente ao prompt original "mais..." \tabularnewline
\hline
\end{tabularx}
\end{center}

\booksubsection{Três casos completos de cenário}

\begin{tcbinfo}{Cenário 1: Assistente pessoal (sugestão de cardápio de café da manhã)}
\textbf{Prompt}:

Você é uma blogueira de café da manhã especializada em nutrição (\textbf{Papel}). Sou uma mãe trabalhadora, com apenas 15 minutos para preparar o café da manhã, tenho um filho de 6 anos e uma sogra de 60 anos em casa (\textbf{Contexto}). Por favor,\textbf{liste 7 combinações de café da manhã diferentes para a semana}(\textbf{Tarefa}). Requisitos para cada café da manhã: (1) tempo de preparo não superior a 15 minutos (2) nutricionalmente equilibrado (proteína + carboidrato + frutas/vegetais) (3) adequado para crianças e idosos (4) ingredientes disponíveis em supermercados comuns (\textbf{Restrições}).

\textbf{Por que funciona}: Papel define o nível profissional, Contexto fornece um ponto em comum para três pessoas, Tarefa especifica claramente o entregável, 4 Restrições garantem que seja executável.

\end{tcbinfo}

\begin{tcbinfo}{Cenário 2: Apoio ao estudo (preparação para o CET-6 de inglês)}
\textbf{Prompt}:

Você é\textbf{um professor experiente de inglês para o CET-6}(Papel). Sou um estudante universitário comum,\textbf{passei no CET-4 com 480 pontos, meta de 500 no CET-6, faltam 3 meses}(Contexto). Por favor,\textbf{elabore um plano de estudos executável de 12 semanas para mim}(Tarefa). Requisitos: (1) definir claramente a distribuição semanal de tempo para listening, speaking, reading e writing (2) indicar a direção principal de foco a cada mês (3) recomendar 3 materiais de estudo adequados (com justificativa em português) (4) apresentar em formato de tabela (Restrições).

\textbf{Por que funciona}: Papel define o contexto profissional, Contexto fornece parâmetros de nota e tempo, Tarefa especifica claramente um “plano faseado”, Restrições exigem “formato de tabela” e “executável”.

\end{tcbinfo}

\begin{tcbinfo}{Cenário 3: Produtividade no trabalho (organização de ata de reunião)}
\textbf{Prompt}:

Você é\textbf{um assistente administrativo sênior}(Papel). Acabei de sair de uma reunião de departamento de 1 hora,\textbf{discutindo o plano de promoção de marketing do Q2}(Contexto).\textbf{Por favor, organize a seguinte gravação da reunião em uma ata estruturada}(Tarefa). Requisitos: (1) dividir em três partes: “Decisões / Itens de ação / Riscos e objeções” (2) nos itens de ação, indicar responsável e prazo (se não mencionado, escrever “a confirmar”) (3) total de palavras não superior a 600 (4) saída em formato markdown (Restrições).

\textbf{Por que funciona}: Papel define a perspectiva administrativa, Contexto dá o tema da reunião, Tarefa especifica claramente a saída estruturada, 4 Restrições tornam o resultado diretamente reutilizável.

\end{tcbinfo}

\textbf{Jogada avançada: transforme a IA em seu “assistente pessoal”}\par

Experimente estes usos avançados para sentir o verdadeiro poder da IA:

\begin{itemize}
  \item \textbf{Otimização por múltiplas rodadas de conversa}: quando não estiver satisfeito com a resposta da IA, diga "Isso está muito técnico, por favor explique de novo de um jeito que uma criança do primário entenda".
  \item \textbf{Cadeia de pensamento (CoT)}: peça para a IA "pensar passo a passo", a qualidade da resposta geralmente melhora: "Por favor, analise este problema de matemática passo a passo".
  \item \textbf{Colaboração entre múltiplas IAs}: use a IA A para escrever o rascunho, a B para polir, a C para traduzir — a combinação traz resultados ainda melhores.
  \item \textbf{Prompt de sistema}: nas configurações da IA, predefina sua identidade, preferências, restrições, etc., para que as respostas se adequem melhor às suas necessidades.
\end{itemize}

\booksection{Limitações e desafios da IA: mantenha o senso crítico}

Embora poderosa, a IA está longe de ser “onipotente”. Reconhecer suas limitações com clareza é essencial para usá-la melhor:

\begin{itemize}
  \item \textbf{Problema de "alucinação"}: a IA pode inventar coisas com convicção (o termo técnico é Hallucination). Ela pode criar eventos históricos inexistentes, citações erradas, estatísticas falsas. Informações importantes devem ser verificadas manualmente.
  \item \textbf{Corte de conhecimento}: a maioria dos modelos de IA tem uma data de corte nos dados de treinamento (como início de 2024). Ela não sabe o que aconteceu depois.
  \item \textbf{Capacidade de raciocínio limitada}: raciocínios complexos com múltiplas etapas, provas matemáticas, armadilhas lógicas — a IA ainda erra com frequência.
  \item \textbf{Riscos de privacidade}: nos serviços de IA na nuvem, suas conversas podem ser usadas para treinar o modelo. Informações sensíveis devem preferencialmente usar IA no dispositivo.
  \item \textbf{Vieses e valores}: a IA pode carregar vieses presentes nos dados de treinamento. É preciso manter o pensamento independente.
\end{itemize}

\begin{tcbinfo}{Dicas de uso de IA}
\begin{itemize}
  \item Conteúdo gerado por IA pode ser impreciso (problema de "alucinação"). Sempre verifique informações importantes.
  \item Não insira informações pessoais privadas na IA (número de documento, senhas, dados bancários, etc.).
  \item Use bem os prompts: quanto mais específica e clara for a descrição, maior a qualidade da resposta da IA.
  \item A IA é uma ferramenta, não uma autoridade. Mantenha o pensamento independente e o senso crítico.
  \item Em decisões cruciais (saúde, questões legais, financeiras), consulte profissionais qualificados em vez de depender da IA.
\end{itemize}

\end{tcbinfo}

\begin{tcbpractice}{Desafio: três exercícios práticos com IA}
\begin{enumerate}
  \item \textbf{Desafio 1: Instrução complexa}: diga ao assistente de voz: "Amanhã às 7h me lembre de levar guarda-chuva, e mude o compromisso das 15h de hoje para as 16h." Veja se ele consegue completar múltiplas operações de uma vez.
  \item \textbf{Desafio 2: Edição de fotos com IA}: tire uma foto de turista com pedestres, use a função "remoção mágica" ou "borracha com IA" para remover os pedestres por completo. Observe o efeito de "preenchimento imaginativo" do fundo pela IA.
  \item \textbf{Desafio 3: Otimização de prompt}: primeiro faça uma pergunta com um prompt simples "escreva um roteiro de viagem", depois use a estrutura CRISPE para escrever um prompt detalhado. Compare a diferença de qualidade entre as duas respostas.
\end{enumerate}

\end{tcbpractice}

\begin{tcbinfo}{Exemplo 2: Fluxo completo de migração de dados ao trocar de celular}
\begin{itemize}
  \item \textbf{Cenário}: iPhone 13 Xiaomi 15, troca de celular entre marcas diferentes.
  \item \textbf{Etapas}: (1) backup no celular antigo via iCloud/WeChat (2) sair do Apple ID + desativar "Buscar" (3) no celular novo, usar o app "Xiaomi Switch" escaneando o QR code (4) selecionar itens a migrar (5) migrar separadamente o histórico de conversas do WeChat com a função "Migração de histórico de conversas" (6) verificar dados críticos (contatos/álbum/WeChat) (7) apagar o celular antigo com criptografia.
  \item \textbf{Resultado esperado}: zero perda de dados, celular antigo pronto para revenda segura.
\end{itemize}

\end{tcbinfo}
